\chapter{Introduction}
\label{ch:native-introduction}

Synthesizing realistic traffic trajectories requires modeling both spatial routing and temporal movement dynamics. In a real-world road network, multiple feasible paths can connect the same origin and destination (OD). Even along a fixed route, vehicles exhibit diverse temporal behaviors—cruising steadily, dwelling at transit stops, accelerating through intersections, or decelerating in response to traffic conditions. Consequently, generating a geometrically connected polyline solves only half of the trajectory synthesis problem. A practically useful generative model must simultaneously respect road-network topology and capture the distinct kinematic distributions and behavioral variations across different traffic participant classes.

The primary output of this work is a synthetic trajectory dataset in the standard \texttt{(time, id, x, y, type)} format, where each identifier corresponds to a complete, continuous journey from origin to destination. In our formulation, the requested participant class governs both road-access permissions and learned movement preferences. Rather than treating synthesis as deterministic point-to-point regression that attempts to recover a single ground-truth future trace, we formulate it as conditional distribution learning: repeated draws for the same request yield diverse, plausible journey realizations. This probabilistic perspective naturally separates our evaluation into distinct dimensions: road-network compliance, physical motion fidelity, within-request diversity, request coverage, and downstream training utility.

This study bridges two active research directions: conditional diffusion models for trajectory generation \cite{ddpm,difftraj,diffrn} and image-based time-series representations for traffic participant recognition \cite{salvadori,ferreira}. Combining these paradigms raises a central research question: \textit{Can spatial and kinematic image representations learned for participant classification improve conditional traffic trajectory synthesis?}

\section{A complete journey with spatial and temporal choices}

Let $c$ denote a traffic participant class, $R$ a permitted directed route on the road network, and $P_R(s)$ its continuous coordinate interpolation at normalized arc-length progress $s\in[0,1]$. We represent a synthesized complete journey of duration $T$ as a timestamped sequence:
\begin{equation}
 \mathcal T=\{(t_i,P_R(s_i),c)\}_{i=0}^{n},\qquad
 t_0=0,\quad t_n=T,\quad
 0=s_0\leq s_1\leq\cdots\leq s_n=1.
 \label{eq:native-trip-definition}
\end{equation}
Across generated samples, the selected route $R$, total travel duration $T$, and intermediate progress profile $\{s_i\}$ vary stochastically. Repeated progress values ($s_i = s_{i-1}$) naturally encode stationary intervals, while varying rates of progress capture acceleration and cruising phases. Because sequence length scales directly with the sampled duration $T$, long-distance transit services are modeled in their entirety rather than truncated to a fixed-length observation window.

To make the generative process transparent and controllable, our framework decomposes trajectory synthesis into four specialized modules: (1) a learned route-preference policy, (2) a training-derived duration prior, (3) a temporal motion diffusion model, and (4) a deterministic road-coordinate decoder. At inference time, a synthesis request specifies only the road network, OD endpoints, participant class, and vehicle access permissions (supplemented by ordered public stop anchors for scheduled buses). First, the route policy samples a valid path from a constructed catalog of road alternatives, and the duration prior samples a total travel time conditioned on route length and class. Next, the diffusion model generates a sequence of progress increments on a uniform $0.25\,\text{s}$ physical-time grid, which the road decoder maps deterministically onto planar coordinates. This modular design ensures strict attribution of model behavior: topological road adherence is guaranteed by catalog construction and decoding, travel-time scaling is governed by the duration prior, and fine-grained speed and stopping dynamics are controlled by the temporal denoiser.

Training on observed route--motion pairs and participant labels enables the denoiser to learn a conditional distribution over complete journeys. At inference time, independent noise draws yield alternative movement realizations for the same request, whose class consistency and kinematic realism are assessed through independent physical metrics and downstream tasks. We develop and evaluate our framework using microscopic traffic simulations on the public Monaco SUMO Traffic (MoST) network \cite{most,sumodocs}, focusing on individual complete journeys conditioned on network attributes and participant classes, with the broader evaluation scope discussed in Chapter~\ref{ch:native-discussion}.

\section{GAF representations for conditional generation}

To investigate visual representation conditioning, we transform each trajectory into two complementary three-channel images comprising Gramian Angular Summation Fields (GASF), Gramian Angular Difference Fields (GADF), and Markov Transition Fields (MTF) \cite{wang,original}. The \textit{spatial view} maps ordered $(x,y)$ coordinates onto a 1D scalar sequence via a space-filling Hilbert curve before computing the image channels. Conversely, the \textit{velocity view} applies the same imaging transforms to a 1D speed profile derived from the trajectory's positions and timestamps. Trained independently to classify the five participant categories, these two CNN encoders capture complementary aspects of vehicle behavior from the same underlying trajectory.

The primary advantage of these image transformations lies in their explicit encoding of pairwise sequence relationships, enabling 2D convolutional encoders to extract global movement patterns that complement a categorical class embedding. Because spatial quantization and run-length compression emphasize route-scale geometry while the velocity view captures the shape and magnitude of whole-trip speed profiles, comparing the two views reveals which aspects of participant-related structure transfer most effectively to motion synthesis.

Nevertheless, discriminative recognition and conditional generation prioritize different properties. A spatial classifier may distinguish buses from motorcycles simply by identifying fixed transit corridors; however, once the route and class are explicitly supplied to the generator, corridor identity offers little additional guidance for synthesizing a realistic within-route speed profile. Moreover, maximizing a classifier's confidence score during sampling can bias generation toward exaggerated, stereotypical class behaviors rather than representative distributions. Our evaluation therefore measures classification accuracy, physical motion fidelity, and downstream data utility as distinct outcomes.

In our main conditioning architecture, we integrate GAF representations directly into the diffusion training objective. Compact 32-dimensional spatial and velocity codes condition every residual block of the temporal denoiser, while the upper stages of the ResNet encoders are jointly adapted through motion reconstruction alongside the participant-classification objective. To support test-time synthesis without access to future target trajectories, we follow the Diffusion Autoencoder paradigm \cite{diffae} and train a conditional latent diffusion prior that samples representation codes given only route geometry, duration, and class inputs.

\section{Simulation data and physical-time representation}

Our experimental corpus is constructed from complete vehicle journeys simulated in SUMO on the public MoST scenario. Leveraging native activity-based demand, class-specific vehicle parameters, and scheduled public transit preserves realistic couplings between participant categories, route choices, and movement dynamics. As shown by our controlled demand comparison in Appendix~\ref{app:source-comparison}, replacing activity-based demand with class-permitted random routes retains local vehicle kinematics but substantially weakens the macroscopic spatial patterns required for participant recognition.

Modeling complete journeys also requires an appropriate temporal representation. Normalizing every trajectory to a fixed sequence length (such as a $257$-point relative-time grid) forces a long bus service onto a much coarser physical time step than a short passenger trip, smoothing out rapid acceleration transients and brief stops prior to learning. Our native-time formulation retains uniform $0.25\,\text{s}$ observations, stationary intervals, and variable sequence lengths. A matched temporal-resolution pilot study evaluates the resulting trade-offs between physical motion fidelity and computational cost.

In total, the expanded learning population comprises $29,895$ complete vehicle trip legs from $24,073$ unique source vehicles across five classes: bus, motorcycle, private car, bicycle, and taxi/ride-hailing. Held-out simulator realizations provide out-of-sample reference trajectories for evaluation under the partitioning rules defined in Chapter 6.

\section{Research questions and contributions}

This thesis addresses two primary research questions:
\begin{enumerate}
    \item Can a modular, road-conditioned diffusion framework synthesize complete multi-class journeys with realistic route preferences and kinematic distributions, and what specific gains arise from preserving native physical time?
    \item Do spatial and velocity GAF representations improve synthesized trajectory quality and downstream data utility beyond standard conditioning on participant class, route geometry, and travel duration?
\end{enumerate}

The main contributions of this work are twofold:
\begin{enumerate}
    \item \textbf{An end-to-end complete-journey synthesis pipeline:} We develop an executable, road-conditioned generative framework combining a permitted road catalog, a learned route-preference policy, a class-conditioned duration prior, variable-length temporal diffusion, and exact road-coordinate decoding. Matched experimental controls cleanly separate structural road compliance, travel-time scaling, and learned temporal motion.
    \item \textbf{A controlled empirical study of GAF conditioning mechanisms:} We systematically evaluate visual trajectory representations across post-hoc candidate selection, frozen-feature supervision, and jointly adapted latent conditioning. Across eight unseen simulation episodes and two model seeds, we benchmark single-view and dual-view GAF conditions against frozen-encoder, raw time-series, shuffled-pairing, and class-centroid controls using multi-timescale physical metrics, qualitative visual inspections, and an independent downstream data-augmentation task.
\end{enumerate}

The remainder of this thesis is organized as follows. Chapter 2 situates our approach within the literature on trajectory diffusion, time-series imaging, and guided generation. Chapter~\ref{ch:native-data} describes the simulation corpus and trajectory representations. Chapters 4 and 5 detail the complete-journey synthesis pipeline and the GAF conditioning modules. Chapter 6 defines the experimental protocol and evaluation metrics, Chapter 7 presents the empirical results, and Chapter 8 discusses the broader implications and scope of our findings.
